\section{Describe the simple SIR epidemic model. Find the condition on which the infection will ultimately die out or spread throughout the population.}
\subsection*{Description of the Simple SIR Epidemic Model}

The \textbf{SIR model} (Kermack and McKendrick, 1927) divides a fixed total population $N$ into three compartmental classes:

\begin{itemize}
    \item $S(t)$: \textbf{Susceptible} individuals who are healthy but vulnerable to catching the disease.
    \item $I(t)$: \textbf{Infected} individuals who currently have the disease and can transmit it to susceptibles.
    \item $R(t)$: \textbf{Removed} (or Recovered) individuals who have either recovered with permanent immunity or died.
\end{itemize}
\noindent\textbf{Model Assumptions:}
\begin{enumerate}
    \item The total population $N$ remains constant over the epidemic time scale (i.e., no births, natural deaths, or migration):
    \[
        S(t) + I(t) + R(t) = N = \text{constant}
    \]
    \item Infection rate is proportional to the rate of contact between susceptible and infected individuals ($\beta S I$).
    \item Recovery rate is directly proportional to the number of infected individuals ($\gamma I$).
\end{enumerate}
\noindent\textbf{Differential Equations:}
\begin{align}
    \frac{dS}{dt} &= -\beta S I \tag{1} \\
    \frac{dI}{dt} &= \beta S I - \gamma I \tag{2} \\
    \frac{dR}{dt} &= \gamma I \tag{3}
\end{align}
\noindent where:
\begin{itemize}
    \item $\beta > 0$ is the \textbf{transmission (infection) rate coefficient}.
    \item $\gamma > 0$ is the \textbf{recovery/removal rate coefficient}.
\end{itemize}
The initial conditions of the above model are:
\[
S(0) = S_0 > 0, \quad I(0) = I_0 > 0, \quad \text{and} \quad R(0) = 0
\]
\noindent Initially, from equation (2):
\begin{align*}
    \frac{dI}{dt} &= \beta S I - \alpha I \\
    \implies \frac{dI}{dt} &= I \beta \left(S - \frac{\alpha}{\beta}\right)
\end{align*}
\noindent Evaluating at $t = 0$:
\begin{align*}
    \left. \frac{dI}{dt} \right|_{t=0} &= I_0 \beta \left(S_0 - \frac{\alpha}{\beta}\right) \tag{iv}
\end{align*}
\noindent From $\frac{dS}{dt} = -\beta S I$, since $\beta, S, I$ are positive, we have $\frac{dS}{dt} < 0$, which implies $S(t)$ is a strictly decreasing function. Therefore:
\[
S(t) < S(0) = S_0 \quad \text{for all } t > 0
\]
\subsection*{Case 1: No Epidemic Outbreak $\left(S_0 \le \frac{\alpha}{\beta}\right)$}
\noindent From equation (iv):
\begin{align*}
    \frac{dI}{dt} = I \beta \left(S - \frac{\alpha}{\beta}\right) &< I \beta \left(S_0 - \frac{\alpha}{\beta}\right) \quad \text{since } S(t) < S_0
\end{align*}
\noindent Now, if $S_0 \le \frac{\alpha}{\beta}$:
\begin{align*}
    \frac{dI}{dt} &< \beta I \left(\frac{\alpha}{\beta} - \frac{\alpha}{\beta}\right) = 0 \\
    \implies \frac{dI}{dt} &< 0 \quad \text{for all } t \ge 0
\end{align*}

\noindent Thus, $I(t)$ is a strictly decreasing function. Integrating gives:
\begin{align*}
    [I]_0^t &< 0 \implies I(t) < I_0 \\
    \implies \lim_{t \to \infty} I(t) &< I_0
\end{align*}
Hence, no epidemic occurs, and the infection dies out.

\vspace{1em}
\subsection*{Case 2: Epidemic Outbreak $\left(S_0 > \frac{\alpha}{\beta}\right)$}

\noindent Let $S_0 > \frac{\alpha}{\beta}$. From equation (iv):
\begin{align*}
    \left. \frac{dI}{dt} \right|_{t=0} &= I_0 \beta \left(S_0 - \frac{\alpha}{\beta}\right) > 0 \\
    \implies \left. \frac{dI}{dt} \right|_{t=0} &> 0
\end{align*}

\noindent Since the derivative is positive at $t = 0$, $I(t)$ is initially an increasing function. Thus, for early times $t > 0$:
\[
I(t) > I(0) = I_0
\]

\noindent The number of infected individuals increases initially. Hence, an **epidemic occurs** and the infection spreads through the population.